\chapter{总结与展望}\label{chap:conclusion}

\section{工作总结}

本文围绕空地协同多功能雷达/电子对抗网络部署优化展开研究。与规则区域内的普通站址选择不同，本文的场景同时包含复杂边界、空地传播差异和覆盖-压制双目标取舍。基于MOPSO-DT框架，论文完成了建模、算法实现、参数调优、多算法对比和边界专项验证。主要结论如下。

（1）复杂区域约束处理。本文把Hertel-Mehlhorn凸分解与垂直交点坐标变换结合使用，将原始非凸区域转化为可编码凸子区域及其内部连续坐标。边界ECR实验显示，该处理相对legacy配置和NSGA-II-DT提高了部分边界任务点覆盖表现；但直接物理空间搜索在单一Boundary ECR上仍取得最高值。因此，本文结论不是“消除边界效应”，而是说明坐标变换能减少非法采样并改善部分边界可达性。

（2）空地异构传播建模。目标评估中区分A2G和G2G链路路径损耗，并兼容雷达方程模型与指数衰减模型。传播模型消融表明，异构建模并不必然提高某个单项指标，它更主要的作用是改变空中节点与地面节点对ECR和$J_{\min}$的贡献分配，避免把两类节点当作同质地面站处理。

（3）MOPSO-DT配置与同预算比较。Challenging场景的同预算对比显示，MOEA/D-DT、SPEA2-DT和NSGA-II-DT在HV、Spacing或非支配解数量上仍有优势；本文配置的价值主要体现在与复杂区域编码、坐标映射和异构传播目标的统一实现，以及相对legacy配置更稳定的前沿分布。

（4）覆盖-压制权衡。Paper-Aligned和Challenging两个场景均观察到ECR与$J_{\min}$的负相关趋势。覆盖率提高通常需要节点更分散，最小压制强度提高则要求节点补强短板任务点。该关系对应电子对抗任务中探测与压制能力的协同分配，可作为当前部署方案选择的依据，但不能外推为所有雷达数量、任务分布和传播条件下的定律。

\section{研究局限性}

本文仍有几方面限制。

\begin{itemize}
    \item 场景维度有限。实验基于二维平面区域，尚未引入真实高程、建筑遮挡和地物类型。三维地形会改变视距关系，也会影响空中节点优势和空地协同平台约束。

    \item 传播模型仍然简化。本文使用自由空间雷达方程或距离指数衰减，没有加入多径、阴影衰落、雨衰和频段相关的随机扰动。在城区或山地环境中，确定性传播模型可能低估局部遮挡。

    \item 随机种子数量不足以支撑严格显著性检验。新增算法对比和边界实验采用3个随机种子，已经比单次运行更稳健，但还不足以给出强统计结论。若用于工程决策，至少需要扩展到更多种子和更多区域形状。

    \item 权衡结论依赖当前任务设置。ECR与$J_{\min}$的负相关来自本文两个主要实验场景。任务点分布、节点数量、功率约束或传播参数变化后，前沿形态可能发生改变。

    \item 部署过程是静态的。本文假设节点部署后固定不动，没有考虑动态对抗环境中的目标移动、平台续航、通信链路重构和在线重规划。

    \item 基线实现偏向统一接口。NSGA-II、MOEA/D和SPEA2均按本文混合变量编码和评估函数实现，便于公平比较，但并不代表这些算法在雷达部署问题上的最强工程调优版本。
\end{itemize}

\section{未来研究方向}

后续工作可以从以下方向继续推进。

\begin{itemize}
    \item 三维与动态部署。把二维部署区域扩展到三维空间，同时优化空中节点水平位置和飞行高度；进一步考虑动态电子对抗环境中的目标移动和在线重部署。

    \item 更细的传播评估。引入数字高程模型、射线追踪或实测信道校准，使A2G/G2G链路参数随地形、频段和遮挡条件变化。

    \item 更充分的统计实验。增加随机种子、区域形状和任务点分布，报告均值、标准差和置信区间，并对主要指标做显著性检验。

    \item 算法组件复用。尝试把MOEA/D或SPEA2的多样性维护机制接入本文的区域分解和坐标变换框架，比较“强优化器+复杂区域映射”的组合效果。

    \item 半实物验证。结合软件无线电平台、信道模拟器或数字孪生环境，检查仿真中得到的部署方案在受控电磁环境下是否保持同样的覆盖和压制趋势。
\end{itemize}
